% Décomposition de l'eau oxygénée : volume de dioxygène dans quatre conditions (exercice 13, chapitre 4)
\begin{tikzpicture}
\begin{axis}[otMesures,otAxesBout,otLecture,width=104mm,height=56mm,xmin=0,xmax=31,ymin=0,ymax=56,
  xtick={0,5,...,30},ytick={0,10,...,50},xlabel={$t$ (min)},ylabel={$V(\mathrm{O_2})$ (mL)}]
\addplot[otGreenDark,line width=1pt,domain=0:30,samples=150] {48*(1-2^(-x/1.0))};
\addplot[otOrange,line width=1pt,domain=0:30,samples=100] {48*(1-2^(-x/8))};
\addplot[otNavy,line width=1pt,domain=0:30,samples=20] {48*(1-2^(-x/1000))};
\node[pcLegende,anchor=south west,text=otGreenDark] at (axis cs:3.2,49.5) {C : catalase};
\node[pcLegende,anchor=north west,text=otOrange!75!black] at (axis cs:15.5,34) {B : ions fer(III)};
\node[pcLegende,anchor=south west,text=otNavy] at (axis cs:11,2.2) {A : sans catalyseur ; D : catalase chauffée};
\end{axis}
\end{tikzpicture}
